\section*{v1.5.0 Projection and Appendix Table Integration Core}
\addcontentsline{toc}{section}{v1.5.0 Projection and Appendix Table Integration Core}

\paragraph{Normalization note.} In Version~\docversion, this block is retained as a \emph{$\mathfrak P$-dominant transport layer}. It introduces no independent generative axioms; rather, it records how closure-licensed structural objects are transported into projection tables, appendix targets, and effective readout interfaces.


The task of v1.5.0 is to convert the earlier derivation-and-appendix
interlock into a more explicit transport layer between the main line and the
technical tables. Earlier versions had already introduced the closure ladder,
the micro-to-macro principle, and the appendix transfer rule. The present stage
adds a stricter rule: every repeated projection, reduction, or read-off step
should be capturable either as a named lemma or as a named table interface with
a declared upstream structural source and a declared downstream effective use.
The point is not stylistic expansion but recoverability. A compressed move must
remain reconstructible.

\subsection*{1. Projection tables as proof interfaces}
Projection tables are not merely descriptive summaries. In v1.5.0 they
receive a precise logical role: a projection table records how a structural
object, already licensed by the closure architecture, may be read as an
effective sector form without pretending that the effective form is primary.
The direction of legitimacy is therefore always
\[
\text{structural source} \Longrightarrow \text{closure-compatible projection}
\Longrightarrow \text{effective sector reading.}
\]
This prevents the later document from treating a Dirac-like, Einstein-like, or
QFT-like formula as an autonomous axiom. The table is only valid if each row
preserves that direction.

\begin{definition}[Projection-table interface]
A projection-table interface is a named record that assigns to a structural
source object its admissible effective readings together with the conditions
under which the passage is allowed, restricted, or still open.
\end{definition}

\begin{lemma}[Projection-interface lemma]
\label{lem:projection-interface-lemma}
Whenever a projection step occurs repeatedly in the main line, the step may be
compressed into a projection-table interface provided that the source object,
closure condition, and interpretation status are all explicitly recorded.
\end{lemma}

\begin{proof}
A repeated projection is proof-relevant only because it preserves a controlled
relation between source and effective reading. If that relation is named once in
a table together with its conditions, the main line may cite the named interface
without losing substance. Compression is legitimate precisely because the table
stores the missing local detail in recoverable form.
\end{proof}

\subsection*{2. Interpretation status discipline}
The earlier reading discipline distinguished between what is structurally
forced, projectively read, interpretively overlaid, and reserved for later proof expansion. In
v1.5.0 this becomes a transport rule for all summary tables and appendix
catalogues.

\begin{proposition}[Interpretation-status transport principle]
\label{prop:interpretation-status-transport}
Every entry transferred from the main line into a table, catalogue, or appendix
record must preserve its interpretation status. A structurally forced item may
not be downgraded into merely interpretive language, and an interpretive or open reading may not
be silently upgraded into a forced statement through tabulation.
\end{proposition}

\begin{proof}
Tables compress syntax, not logical rank. If the interpretation status were
allowed to drift during compression, then the appendix would alter rather than
preserve the proof surface. The transfer must therefore carry both content and
status labels together.
\end{proof}

\paragraph{Practical consequence.}
Projection tables must visibly mark at least three columns: structural source,
allowed effective reading, and status label. Optional columns may record frame
conditions, probe restrictions, or unresolved hypotheses, but the status label
is mandatory.

\subsection*{3. Appendix-citation lattice}
The appendix now receives a more explicit citation lattice. Instead of treating
the appendix as a linear storage area, v1.5.0 organizes it as a small
network of recoverable proof targets:
\begin{enumerate}[leftmargin=2em]
  \item \textbf{lemma targets} for local derivation moves;
  \item \textbf{tool targets} for normalized rechen-tricks;
  \item \textbf{table targets} for repeated projection or reduction patterns;
  \item \textbf{schematic targets} for probe geometry and search-space layouts.
\end{enumerate}
A compressed sentence in the main text is acceptable only when it points to one
of these four target types.

\begin{lemma}[Appendix-citation lattice lemma]
\label{lem:appendix-citation-lattice-lemma}
A compressed main-text step is proof-disciplined if and only if it can be
expanded along at least one explicit appendix citation path ending in a lemma,
tool, table, or schematic target.
\end{lemma}

\begin{proof}
The project rule forbids concealment by relocation. Therefore every compressed
step must have a recoverable destination. The four target types above exhaust
the relevant forms of local expansion used in the present monolith. If no such
path exists, the compression is not yet justified.
\end{proof}

\subsection*{4. Projection robustness across frames and probes}
A repeated source of ambiguity is the temptation to read the same effective term
as though it had identical status in all frames and in the probe sector. The
present release therefore states a robustness rule.

\begin{proposition}[Projection-robustness principle]
\label{prop:projection-robustness}
An effective reading is robust only if it survives both frame-shift transport and
probe-sector restriction without leaving the closure-compatible class.
\end{proposition}

\begin{proof}
A term read in one frame but destroyed by admissible frame transport cannot be a
stable structural projection. Likewise, a reading that collapses once the probe
sector imposes orthogonality and closure tests is not globally licensed. Robust
readings are exactly those that persist under both operations.
\end{proof}

\paragraph{Use in later versions.}
This principle is meant to control later expansions of matter, gauge, mass, and
Higgs readings. The question is no longer only whether a reading is suggestive,
but whether it survives transport across the disciplined structural operations
already fixed by the framework.

\subsection*{5. v1.5.0 readiness criteria}
This release should be regarded as complete in working form when the following
conditions hold simultaneously:
\begin{enumerate}[leftmargin=2em]
  \item the version stamp and DOI are updated to v1.5.0;
  \item projection tables are assigned an explicit proof-interface role;
  \item interpretation status is required to survive table and appendix transfer;
  \item the appendix-citation lattice is stated as a reusable compression rule;
  \item projection robustness is tied to both frame transport and probe tests;
  \item the new material is appended without shortening the inherited monolith.
\end{enumerate}

These conditions are now satisfied in the present monolith. Relative to v1.4.0,
v1.5.0 marks the point where appendix tables and projection summaries are
integrated into the proof discipline itself rather than serving only as support
material.

\clearpage
\clearpage


\clearpage
